% Generated draft: requires mathematical and source review.
\documentclass[11pt]{article}
\usepackage[margin=25mm]{geometry}
\usepackage{fontspec}
\setmainfont{Latin Modern Roman}
\usepackage{amsmath,amssymb,mathtools,unicode-math}
\setmathfont{Latin Modern Math}
\usepackage{enumitem,needspace,xurl,hyperref}
\setlength{\parindent}{0pt}
\setlength{\parskip}{5pt}
\setlength{\emergencystretch}{3em}
\newcommand{\N}{\mathbb{N}}
\newcommand{\Z}{\mathbb{Z}}
\newcommand{\Q}{\mathbb{Q}}
\newcommand{\R}{\mathbb{R}}
\newcommand{\C}{\mathbb{C}}
\newcommand{\F}{\mathbb{F}}
\newcommand{\A}{\mathbb{A}}
\newcommand{\PP}{\mathbb{P}}
\newcommand{\E}{\mathbb{E}}
\newcommand{\eps}{\varepsilon}
\newcommand{\dd}{\,\mathrm{d}}
\begin{document}
\section*{UnsolvedMath ID: 30002193}
\textbf{Source problem code:} catalogue entry
\subsection{Definitions and mathematical statement}
\textbf{Interpretive conventions.} The supplied record does not specify the ground field, analytic category, deformation equivalence, or resolution convention. The following is a conditional formulation in the customary setting of complex analytic surface germs, using minimal good resolutions and analytic isomorphism. These conventions require confirmation against the cited sources.

Let $(X,x)$ be a singular normal complex analytic germ of dimension two. Call $(X,x)$ \emph{simple} if, in a sufficiently small representative of a semiuniversal deformation, the germs at points of nearby fibres represent only finitely many analytic isomorphism classes. Call it \emph{rigid} if every sufficiently small deformation is locally analytically trivial. These definitions concern deformations of the singularity, without imposing an additional condition such as preservation of the resolution graph.

A normal surface singularity is \emph{rational} if, for a resolution $\pi:\widetilde X\to X$, one has $R^1\pi_*\mathcal O_{\widetilde X}=0$ near $x$. In this formulation, a \emph{rational double point} is a rational surface singularity of multiplicity two, and a \emph{rational triple point} is a rational surface singularity of multiplicity three.

For a minimal good resolution, let $\Gamma(X,x)$ be its weighted dual graph. Its vertices correspond to irreducible exceptional curves $E_v$, its edges record intersections, and each vertex has self-intersection weight $w_v=E_v^2$ and genus label. Say that a graph is obtained from another by \emph{decreasing vertex weights} if their vertices, edges, and genus labels can be identified and their weights satisfy
\[
 w'_v=w_v-k_v,\qquad k_v\in\Z_{\geq 0}
 \quad\text{for every vertex }v.
\]
This permits all $k_v$ to vanish, so the original graphs are included; the supplied phrase ``making some vertex weights more negative'' does not explicitly settle this convention.

\textbf{Proposed classification, conditional on these conventions.} For every singular normal complex analytic surface germ $(X,x)$,
\[
\begin{aligned}
 (X,x)\text{ is simple}
 \quad\Longleftrightarrow\quad& (X,x)\text{ is rational, and there exists a rational}\\
 &\text{double point or rational triple point }(Y,y)\\
 &\text{such that }\Gamma(X,x)\text{ is obtained from }\Gamma(Y,y)\\
 &\text{by decreasing vertex weights as defined above.}
\end{aligned}
\]
The supplied statement also asserts, as an intended consequence, that no singular normal surface germ is rigid:
\[
 \forall (X,x)\text{ in the stated category},\qquad
 (X,x)\text{ is not rigid}.
\]
Here the restriction to singular germs is essential: smooth germs are excluded. The supplied record does not provide the argument relating the classification to this additional assertion. Both assertions are presented as conjectural in the supplied dated assessment, not as established theorems.
\subsection{Short English statement}
The conjecture classifies simple normal surface singularities as precisely the rational ones whose resolution graphs come from rational double or triple points by decreasing self-intersection weights. It also asserts the nonexistence of rigid singular normal surface germs. The complex analytic setting, deformation definitions, resolution convention, and inclusion of unchanged graphs used here require source confirmation.
\subsection{Sources}
\begin{itemize}
\item Local UnsolvedMath registry, permanent ID \texttt{30002193}, title \emph{Classification of Simple Normal Surface Singularities}, source problem code \texttt{catalogue entry}, in the \emph{Oberwolfach Reports Open Problems} collection. The supplied record lists proposed year 2012 and names no proposer.
\item Supplied report citation: \emph{Singularities (2013)}, Oberwolfach Reports, DOI \texttt{10.4171/owr/2012/46}, pp.~2799--2866. Supplied original-report link: \url{https://doi.org/10.4171/owr/2012/46}. The citation's title year and the record's proposed year differ; both are retained as supplied.
\item Supplied literature reference: Stevens, \emph{Simple surface singularities}, \url{https://arxiv.org/abs/1303.0692}.
\item Supplied literature assessment dated \texttt{2026-08-21}: the classification and the related nonexistence assertion remain conjectural; Stevens proved one direction generally and special classes in the other direction. The record does not identify these directions or special classes. This is supplied dated curation, not an independently verified literature assessment. The background explicitly requests independent verification of time-sensitive status and cited claims.
\item The registry labels the problem \texttt{open} and its statement \texttt{corrected\_verified}, with a supplied claim of manual semantic, typographic, and LaTeX review. These are source metadata and do not establish mathematical correctness or independently verified current status. No external source was visited for this draft.
\end{itemize}
\end{document}
